\documentclass[a4paper, UTF-8]{article}
\usepackage[margin=1in]{geometry}
\usepackage{algorithm}
\usepackage{algpseudocode}
\usepackage{amsmath}
\usepackage{amsfonts}
\usepackage{amsthm}

\usepackage[utf8]{inputenc} % allow utf-8 input
\usepackage[T1]{fontenc}    % use 8-bit T1 fonts
\usepackage{hyperref}       % hyperlinks
\usepackage{url}            % simple URL typesetting
\usepackage{booktabs}       % professional-quality tables
\usepackage{amsfonts}       % blackboard math symbols
\usepackage{nicefrac}       % compact symbols for 1/2, etc.
\usepackage{microtype}      % microtypography
\usepackage{xcolor}         % colors

\theoremstyle{plain}
\newtheorem{theorem}{\indent Theorem}[section]
\newtheorem{lemma}[theorem]{\indent Lemma}
\newtheorem{assumption}{\indent Assumption}[section]
\newtheorem{note}[theorem]{\indent Notation}
\newtheorem{proposition}[theorem]{\indent Proposition}
\newtheorem{corollary}[theorem]{\indent Corollary}
\newtheorem{definition}{\indent Definition}[section]
\newtheorem{example}{\indent Example}[section]
\newtheorem{remark}{\indent Remark}[section]
\newenvironment{solution}{\begin{proof}[\indent\bf Solution]}{\end{proof}}
\renewcommand{\proofname}{\indent\bf Proof}

\begin{document}

\title{Assignment 1}

\author{Yanqiao Chen\\ 
  Department of Computer Science and Engineering\\
  Southern University of Science and Technology\\
  \texttt{12412115@mail.sustech.edu.cn}
  }


\maketitle

\section{Page 27, Problem 2}

\subsection{a}

\begin{align*}
\Omega = &\{(1,1), (1,2), (1,3), (1,4), (1,5), (1,6), (2,1), (2,2), (2,3), (2,4), (2,5),(2,6),\\ 
         & (3,1), (3,2), (3,3), (3,4), (3,5), (3,6), (4,1), (4,2), (4,3), (4,4), (4,5), (4,6),\\ 
         & (5,1), (5,2), (5,3), (5,4), (5,5), (5,6), (6,1), (6,2), (6,3), (6,4), (6,5), (6,6)\}
\end{align*}

\subsection{b}

\begin{enumerate}
  \item \begin{align*}A = &\{(1,4), (1,5), (1,6), (2,3), (2,4), (2,5), (2,6), (3,4), (3,5), (3,6), (4, 1), (4, 2), (4, 3), (4,4), (4,5),\\ & (4,6), (5, 1), (5, 2), (5, 3), (5,4), (5,5), (5,6), (6, 1), (6, 2), (6, 3), (6,4), (6,5), (6,6)\} \end{align*}
  \item $B = \{(2,1), (3,2), (3,1), (4,3), (4,2), (4,1), (5,4), (5,3), (5,2), (5,1), (6,5), (6,4), (6,3), (6,2), (6,1)\}$
  \item $C = \{(4,1), (4,2), (4,3), (4,4), (4,5), (4,6)\}$
\end{enumerate}


\subsection{c}
\begin{enumerate}
  \item  $A\cap C = \{(4,4), (4,5), (4,6), (4,1), (4,2), (4,3)\}$
  \item \begin{align*}B \cup C = &\{(2,1), (3,2), (3,1), (4,3), (4,2), (4,1), (5,4), (5,3), (5,2), (5,1), (6,5), (6,4), (6,3), (6,2),\\ & (6,1), (4,4), (4,5), (4,6)\}\end{align*}
  \item \begin{align*}A \cap (B\cup C) = &\{(3,2), (4,1), (4,2), (4,3), (4,4), (4,5), (4,6), (5,1), (5,2), (5,3), (5,4), (6,1),\\ & (6,2), (6,3), (6,4), (6,5)\}\end{align*}
\end{enumerate}

\section{Page 27, Problem 5}

$$
C = (A \cup B) \cap (A \cap B)^c
$$

\section{Supplementary Problem 1}

$$
\overline{A \cup B} = \overline{A} \cap \overline{B} = A \cap B
$$

Then 
$$
\overline{A \cup B} = A \cap B
$$

Since
$$
\overline{A \cup B} \subseteq \overline{A}
$$
and
$$
A \cap B \subseteq A
$$
and
$$
 A\cap \overline{A} = \emptyset
$$

Then as the bigger set $A$ and $\overline{A}$ have no intersection, the smaller set $A \cap B$ and $\overline{A \cup B}$ also have no intersection. But they 
are the same set, so they must be empty. Formally, $(A \cap B) \cap \overline{A \cup B} = (A \cap B) \cap (A \cap B) = \emptyset$, 
we conclude that
$$
A \cap B = \emptyset
$$
and
$$
\overline{A \cup B} = \emptyset
$$

Thus 
$$
A \cup B = \Omega
$$

\section{Supplementary Problem 2}

Let

\begin{align*} 
  & B_1 = A_1 \\
  & B_2 = A_2 \cap A_1^c \\
  & B_3 = A_3 \cap A_1^c \cap A_2^c \\
  & \cdots \\
  & B_n = A_n \cap \bigcap_{i=1}^{n-1} A_i^c
\end{align*}

Then we can verify that $B_i$'s are disjoint and $\bigcup_{i=1}^n A_i = \bigcup_{i=1}^n B_i$. So we have

$$
A_1 \cup A_2 \cup \cdots \cup A_n = B_1 \cup B_2 \cup \cdots \cup B_n
$$

\end{document}